\section{Engineering Interpretation}
\label{sec:limits}
The evidence addresses the three distinctions motivating the design.
Profile tests exercise production-scoped authority; admission ablation
isolates the pre-write interval from readback. The publish-only challenge
separates transport acceptance from endpoint confirmation, and paired
trajectories link interventions to observations and retained recipes.
These results support an inspectable execution contract on the exercised
routes.

An LLM, scripted policy, or deterministic controller supplies a candidate
through its authorized interface; the platform applies the execution
contract. Process reuse changes the signal map, objective, and guards while
retaining production context and execution records. In Q6, observed
acceptance of approved trials precedes retention of the recipe version.

Human participation is represented through explicit authority and
attributed decisions. This supports inspection of what an approver
authorized and what evidence the proposer later used. Automation research
distinguishes such workflow structure from operator comprehension,
complacency, and allocation of responsibility
\cite{parasuramanmanzey2010complacency,parasuraman2000model,bainbridge1983ironies}.
The server-side boundary resolves execution authority independently of
prompts and retrieved documents \cite{greshake2023injection,liu2024pinjection}.
Deployment combines this boundary with credential management, endpoint
security, and IT/OT separation \cite{iec62443-3-3}. The supervisory layer
coexists with independently engineered PLC interlocks and safety
instrumented systems; IEC~61511 governs the latter's process-safety
lifecycle \cite{iec61511}.

\noindent\textit{Experimental scope.}
The process evidence uses accelerated simulators, scripted approvals, and
sequential trajectories. Q1--Q4 cover the historical direct-node interface;
current profile and composition tests cover software behavior; Q6 covers
one recipe trajectory. No single current-version end-to-end study jointly
varies strategy, protocol, and production context, and a complete current
AML-exploration control--response trial is not archived. Q6's observed
guards, snapshots, and cadence settings
are specified in Section~\ref{sec:recipeexperiment}. These evidence units
establish the exercised
mechanisms without estimating comparative optimization benefit, fault
frequency, human performance, hard real-time behavior, or physical transfer.
Matched initial states, repeated independent runs, and a fixed-policy or
SCADA comparator would support those further estimates. Established process
benchmarks such as Tennessee Eastman \cite{downs1993tep} provide an
additional setting for that evaluation.

\section{Conclusion}
\label{sec:conclusion}
The proposed framework connects replaceable industrial agent strategies
to production-scoped authority, graded command evidence, and distinct
node/recipe records. Admission and confirmation experiments demonstrate the
separation of pre-write permission from post-write evidence. Controller and
process missions exercise strategy and configuration reuse, while the
injection-molding trajectory links approved trials to an accepted recipe.
Together, these simulator-backed results support a common execution and
evidence contract for inspectable industrial agent supervision.

\noindent\textit{Data availability and AI assistance.}
Run archives, original Q2 engine logs, replay entries, configuration notes,
and claim-to-source mappings accompany the separate anonymous experimental
evidence bundle. The manuscript source package contains plotting sources
and figure mappings. An OpenAI language model assisted
with manuscript organization, language revision, and plotting-code
preparation. OpenAI image generation drafted the compositions of
Figs.~\ref{fig:arch} and~\ref{fig:runtime}; the final artwork was redrawn
as editable Python-generated vectors and checked against the implementation.
The authors remain responsible for the final submission.
